\section{Angels, Templars, and Joseph de Maistre}

\begin{quotex}
We must be ready for a profound event in the divine order, toward which we are marching with an accelerated speed that must strike all observers. Terrible oracles already announce that \emph{the time has come}.
\flright{\textsc{Joseph de Maistre}, \emph{Soirées de Saint-Petersbourg}}
\end{quotex}


Herein, we review the following books:

\begin{itemize}


\item Seven Archangels: Rhythms of Inspiration in the History of Culture and Nature\footnote{\url{http://www.sophia.sk/en/kniha/sedem-archanjelov}} by \textbf{Emil Pales}

184 pages in full colour, linen hardcover with dust jacket, format 190 x 135 mm, 115 colour illustrations, 105 diagrams

\item Contra Mundum: Joseph de Maistre \& The Birth of Tradition\footnote{\url{https://angelicopress.org/product/contra-mundum/}} by \textbf{Thomas Garrett Isham}

\item The Templars: Knights of Christ\footnote{\url{https://angeluspress.org/products/templars-knights-christ}} by Regine Pernoud

\item \emph{Studies in Freemasonry \& the Compagnonnage} by \textbf{Rene Guenon}
\end{itemize}


\paragraph{Angelology}


First of all, we should mention that the book itself is a work of art. A finely crafted hardcover printed on glossy paper with dozens of illustrations and diagrams. It even comes with a purple ribbon for bookmarking. The appendix is also of interest. It is comprised of a series of seven images of the angels by the artist David Newbatt, along with a detailed interpretation.

Emil Pales is a philologist and mathematician living in Slovakia. His aim is to continue the theology of Sophia which was developed by Vladimir Solovyov. This involves, as we have pointed out, the synthesis of science, art, philosophy, and religion. He claims Goethe, the biologist Adolph Portmann, and Rudolf Steiner, as other influences.

The main idea is that the seven named Archangels are really the spirits of the seven traditional planets, viz., the Greek gods associated with them; he even traces them back to the Babylonians. Each of the Archangels rules or influences a specific time period in human history. Using statistical techniques, Dr. Pales has documented periodic cross-cultural changes in human history. This book is an abridgment of the 1500 page Slovak version, which has more extensive supporting documentation. We miss the days when the educated would communicate and publish in Latin; there are only so many languages one can learn.

Dr. Pales offers this explanation:

\begin{quotex}
Angelology is a spiritual astronomy. It is a science about introspective observation of phenomena in the spiritual sky.
\end{quotex}


Star and angel are used interchangeably in the bible, since angels are the spiritual intelligence of the stars. Just as the planets and stars have their laws and cycles, so too does the activity of the angels. Their activity can be noticed on three distinct levels:

\begin{enumerate}


\item The evolution of the universe

\item Changes in human culture

\item The developmental stages in the life of the human being
\end{enumerate}


The correspondence between the planets and the angelic intelligences was pointed out in\emph{Tantra Vidya\footnote{\url{https://www.gornahoor.net/?p=8618}} by \textbf{Oscar Hinze}}. Using a totally different approach, he demonstrated that the correspondences between the chakras and the planets could not have arisen by chance. This would account for the third level. Dr. Pales goes further, showing that levels one and two are not random processes, but are rather the manifestation of deeper metaphysical laws. For reference, Table \ref{tab:PlanetCorrespondences} has the correspondences noted in the book.

\begin{longtable}[]{@{}cccp{.5\textwidth}@{}}
\caption{Planet correspondences}
\label{tab:PlanetCorrespondences}\\
\toprule
\textbf{Planet} & \textbf{Babylon} & \textbf{Angel} & \textbf{Description} \\
\midrule
\endhead

Moon & Sin & Gabriel & The archangel of the Moon - Enlivener of nature - Inspirer of childhood and the Baroque - The most picturesque of the angels - Unrecognized by scientists and philosophers - The queen of the heavens is alive \\

Mercury & Nabu & Raphael & The archangel of Mercury - The patron of physicians, botanists and linguists - The geological age of Mercury. The birth of Gothic - The inspirer of rationalism and the Enlightenment - The joker and protector of youth \\

Venus & Innana & Anael & The archetype of love and beauty - The spirit of puberty - The muse of fine arts. Soul copper - Compassionate saints - Romantics, mystics and revolutionaries - Patron of alchemists \\

Sun & Shamash & Michael & The sun of truth - The spirit of the Paleozoic era - Michaelic art - Sun cults - The inspirer of philosophy - The history of democracy - The task of the present day \\

Mars & Nergal & Samael & Gods of war - The geological age of dragons - The universal conflagration - The rhythm of the goddess Eris - The twentieth century - New knighthood  \\

Jupiter & Marduk & Zachariel & The god of the celestial dome - The geological age of Sagittarius - The patron of Indo-Europeans - The great kings - Renaissance aesthetics and thinking - The inspirer of geometry and geography - The soul blue of those in their fifties \\

Saturn & Ningirsu & Oriphiel & Soul lead - The geological age of Capricorn - The patron of hermits and emperors - Oriphielic geometry: the circle - Saturnian sciences - Every man is the architect of his own fortunes \\

\bottomrule
\end{longtable}

For the modern mind, this is all absurdity. Dr. Pales points out the three stages of decline that leads to our incredulity:

\begin{itemize}


\item First, direct experience of angels

\item Then, mere verbal assent to their existence

\item Finally, disbelief
\end{itemize}


Like Rene Guenon, Dr. Pales, traces the devolution of our understanding of angels to the waning of the Middle Ages. Apart from a few, people believed in angels, but without any gnosis. Angels then become the subjects of schlock art, but without any relevance to natural processes and world affairs. For example, I recently listened to a youtube talk by a Catholic doctor of theology. He scoffed at the notion that there are cosmic forces at play in the world.

To recover an interest in angels, in is imperative to understand who and what they are. Viz., they are discarnate intelligences. That means, we experience them in our thoughts, moods, inner images, etc. These inner experiences are so close to our identities that we seldom take care to obverse their real source. This book will serve as a guide to become aware of angelic influences in cosmic evolution, human history,

For the still skeptical, there are the many graphs and histograms demonstrating the thesis of the book by using contemporary research tools. Nevertheless, the events described are not due to impersonal, mechanical causes. Since angels have intelligence and will, there is ipso facto a moral cause behind the world process. This leads us into a discussion of Joseph de Maistre.

\paragraph{Joseph de Maistre}


\textbf{Thomas Garett Isham} wrote the book \emph{Contra Mundum} about \textbf{Joseph de Maistre} and his ties to esoteric teaching. In particular, he compares Maistre to Guenon in a fruitful way. He has previously written a few books on the Enneagram.

Thus, in this book we get an esoteric perspective on Maistre. The most influential essays on Maistre have been by the atheist Isaiah Berlin and by the nihilist Emil Cioran. As Mr. Isham points out, they both see in Maistre only a proto-fascist; yet Berlin condemns him and Cioran praises him. Obviously, such an anachronism does a disservice to Maistre. Berlin and Cioran are unaware of the Maistre's deeper spiritual project.

First of all, he was intransigently opposed to the philosophes of the Enlightenment. Although a conventionally religious Catholic, he recognized the importance of the esoteric dimension, which was not opposed to exoteric practice. Thus for years, he was a member of a Masonic lodge. But that was at a time when it was more focused on esoteric ideas, before the Masons became a secretive political power. In particular, Maistre was influenced by the illuminist writings of Louis Claude de Saint-Martin. As Mr. Isham pointed out:

\begin{quotex}
members were drawn for weightier reasons and they found the churches spiritually mediocre and ill-suited to the mentality of the age.
\end{quotex}


Maistre proposed three degrees on Masonry:

\begin{itemize}


\item The promotion of benevolent acts and the study of morals and politics.

\item The promotion of the reunion of the Christian churches and sects.

\item Seeking the revelation of revelation, the sublime knowledge that constitutes knowledge of man at the highest level, a transcendent Christianity based on factual research and metaphysical knowledge.
\end{itemize}


Maistre further adds:

\begin{quotex}
This Science does not differ essentially from ancient Greek or Egyptian initiation.
\end{quotex}


Mr. Isham writes about Maistre's interest in ancient pagan and Christian antiquity, particularly his high regard for Origen. He writes:

\begin{quotex}
Maistre shared Origen's belief that the planets were ``alive,'' part of a cosmos ``created by and for intelligence.'' Imbued with such notions Maistre went so far as to urge the Platonic idea that the stars themselves were divinities and to hold that at least some early Christians found this belief conformable to dogma and of use in their initiations.
\end{quotex}


Throughout his book, Mr. Isham points out some ``amazing'' parallels between Maistre and Rene Guenon. Obviously, they had common interests: the Primordial Tradition, initiatic organizations, esoterism, Masonry, spiritual hierarchy, reaction against Modernity and the Enlightenment. He then adds an extensive review of Guenon's Spiritual Authority and Temporal Power. What Guenon makes explicit in a systematic way, Maistre held implicitly, scattered throughout his writings.

In our time, constitutions, whose origins are public and allegedly clear, are believed to be the foundations of the political system. Yet that is delusional, since there is debate about what precisely the constitution means. Hence, there is something deeper, and not fully understood, that is the real foundation. Mr. Isham explains Maistre's opposed view:

\begin{quotex}
In Maistre's view, the origins of both spiritual and temporal authority are necessarily shrouded in mystery. For obscure origins are the hallmark of legitimacy. Miracle, ritual, the favor of heaven; all are part and parcel of the traditional order. Hence legendary or mythical beginnings, he insists, more adequately ensure the wisdom of the governors and the consent of the governed than the shaky parliamentary systems and man-made constitutions of modern dreamers.
\end{quotex}


In \emph{Contra Mundum}, we have a work that deals with Maistre's ideas as he may have experienced them, rather than as an academic exercise.

\paragraph{The Templars}


The story of the Knights of the Temple of Solomon has had many variations over the years. Some consider them as a holy Order of Knights that protected pilgrims to the Holy Land. Others see in them the founders of usury and deviant spiritual practices. More outlandish are the conspiracy theories described in various fictional accounts.

Curiously, the ultra-traditional SSPX organization has promoted the recovery of the Templars by making available a CD and a book on them. The book \emph{The Templars: Knights of Christ} by Regine Pernoud provides an historical record of the Templars, based on primary documentation. As such, it is quite matter-of-fact in its approach. Although such an approach is informative, and even necessary, Pernoud does not address the deeper spiritual nature of the Templars. Presumably that is because the Templars did not commit such topics to the written word. Moreover, in an initiatic organization, the leaders would have been of a higher grade. The neophyte knights, on the other hand, would not have been privy to their deeper knowledge.

There is extensive documentation about the initiation rites, the daily life, the architecture of the Templar buildings, their many battles, their banking activities, and finally their arrest and trials. There is too much detail, e.g., their daily food productions, to mention here. What is amazing is the sheer number of Templar monasteries scattered across Europe. Their unique life as both monks and warriors must have appealed to a large number of men. In addition to the knights, a large supporting staff of sergeants, squires, and sundry other workers were necessary to sustain them. Although they were the recipients of largesse from the wealthy, each monastery was self-sustaining.

The story of their decline is disheartening. After years of success in the Holy Land, they started losing battles. Rivalry with the Hospitallers led to unnecessary battles. King Philip the Fair coveted their financial wealth. He had the Templars in France arrested on trumped up charges, with confessions made under torture. In that way, he aimed to confiscate their wealth. He pressured Pope Clement V to disband the order, thereby inverting the proper order between the spiritual and temporal authorities. Thomas Isham makes this clear:

\begin{quotex}
By these actions, temporal power began to use spiritual authority for its own ends. Guenon says it is instructive that Dante attributed Philip's action to cupidity, the vice of greed, the vice not of a Kshatriya but of a Vaishya. ``We could say that when they enter a state of revolt, the Kshatriyas, as it were, degrade themselves, losing their own character and taking on that of a lower caste.
\end{quotex}


\paragraph{Templars and the Esoteric Tradition}


In \emph{Studies in Freemasonry \& the Compagnonnage}, Rene Guenon includes a chapter on Joseph de Maistre. For our purposes, we are more interested in what Guenon has to say about the Templars and their relationship to Masonry. He quotes Maistre: ``What does the destruction of the Order of the Templars matter to the Universe?'' Guenon then responds:

\begin{quotex}
But it matters very much, since this marks the point at which the West broke with its own initiatic tradition reaches back, a rupture that is truly the primary cause of the intellectual deviation of the modern world. This deviation in fact goes back further than the Renaissance, which marks only one of the principal stages, the fourteenth century being the actual starting point.
\end{quotex}


This corresponds to the time frame mention by Emil Pales. Although Maistre maintained their existence, he knew nothing of the actual means of the transmission of initiatic doctrines nor of the representatives of the true spiritual hierarchy. In the case of the Masons, Maistre noted that the beginning of Masonry coincided with the destruction of the Templars, but did not really connect those events.

The irony is that the last Grandmaster of the Templars, Jacques de Molay, who was executed by the Church, then became the spiritual father of the Freemasons, the sworn enemy of the Church. Moreover, both Masonry and the Church seem to have lost their initiatic character. Nevertheless, Guenon rather optimistically believed that the Great Work could be resumed in some other form. He explained:

\begin{quotex}
Is this to say that such a Plan will never again be undertaken in one or another form by some organization with a truly initiatic character and possessing an Ariadne's Thread to guide it through the labyrinth of innumerable forms under which the sole Tradition is hidden, in order finally to recover the Lost Word and leave the Light of the Shadows, the Order of Chaos?
\end{quotex}



\flrightit{Posted on 2019-08-11 by Cologero}

\begin{center}* * *\end{center}

\begin{footnotesize}
\begin{sffamily}
\texttt{admin@warofart.xyz on 2019-08-12 at 00:11 said:}

This is honestly making me consider giving Masonry a shot. Anyone have any real experience with it? Is it worth the trouble? (Also reading Mouravieff, but I'm still trying to get the sages pose consistent, and it's slow going.)

\hfill

\texttt{Cologero on 2019-08-12 at 10:39 said:}

I doubt you'll find what you're looking for in contemporary Masonry.

As we pointed out: Guenon rather optimistically believed that the Great Work could be resumed in some other form

That new form should be attached to a genuine Tradition. On this blog, we've been trying to outline what that Great Work might look like.

\hfill

\texttt{Masonguy on 2019-08-14 at 03:24 said:}

If you find what they call a ``traditional'' lodge. There was an attempt made recently at my lodge to make it more traditional which I supported, but it was squelched by the majority who view esotericism with suspicion. Go figure. There are all kinds of lodges just like all kinds of people, and frankly you have to take it on a lodge by lodge basis. Some are very into the esoteric, while others are little more than old boy cigar smoking clubs or just dining clubs where the guys get to get away from the wives at least once a month.

\hfill

\texttt{Fallen adam on 2019-08-15 at 09:08 said:}

I've also been having problems with achieving the posture of the sage. I use a hard wood stool og about 35 cm, but the problem is that if I can't find a comfortable and relaxed position while having my legs crossed. Should I keep trying, or just do with my legs having a 90 degrees angle at the knee, which I've found to be easier to do?

\hfill

\texttt{Sviatoslav on 2019-08-16 at 09:28 said:}

Great to hear that my favorite website Gornahoor is familiar with Emil Pales. Here in Slovakia, he even has a school of Angelology and his books and podcasts deeply influenced me. He`s a profound person.

Although some healthy masonry lodges probably exist and that can be an option for inner growth of a few individuals, freemasonry as a whole cannot become a vehicle for a Traditional message. It`s too corrupt and essentially nonesoteric.

Purely esoteric order in modern circumstances is very hard to establish and maintain, but I think that restoration of warrior-initiatic order as Templars were, isn`t so unrealistic even nowadays. Such an organization could be founded with few idealistic men who are as well capable in practical areas -- making money, engineering, management, military, marketing. So the organization would be self-funded and would operate in a small, but effective team. And then is possible to create more exoteric organizations, e.g. metapolitical, political and through them to influence society.

I think this would be a task for men of warrior nature with ideological help of priests. Unexpectedly, globalization could help us with international recruitment of members (foremost Europeans countries). Unlike all modern mass organizations, which are dependent on society conditions a small elitist group can operate even among ruins.

Has someone opinion on this?

\hfill

\texttt{Sviatoslav on 2019-08-17 at 17:55 said:}

I just want to add to my previous comment, that this website \url{https://www.bernardoption.com} is trying to do the same thing.

After Gornahoor another excellent work. Hopefully, I will find more valuable information to study. I have found a few interesting ideas here \url{https://www.knightstemplarorder.org/} but also a lot of dubiously egalitarian views (membership for women, muslims and so).

\hfill

\texttt{Sviatoslav on 2019-08-21 at 08:16 said:}

That could be done. I will try to upload on youtube an English version of his best podcast.

\hfill

\texttt{Cologero on 2019-08-21 at 09:17 said:}

Keep in mind that you can hear Emil Pales live on Monday, August 26... and even ask questions.

Much better than a podcast.

\hfill

\texttt{Cologero on 2019-08-21 at 22:01 said:}

Make yourself familiar with the Angels, and behold them frequently in spirit. without being seen, they are present with you. St. Francis de Sales

\hfill

\texttt{Boreas on 2019-09-27 at 03:44 said:}

I would love to hear your take on de Maistre's Ultramontanism.

\hfill

\texttt{John Fitzgerald on 2021-03-08 at 15:27 said:}

The idea of the stars as divinities (De Maistre/Origen) feels congruent and right to me. CS Lewis articulated something very close to this in The Last Battle when Aslan calls time on the Narnian world and beckons the stars home --

`Stars began falling all around them. But stars in that world are not the great flaming globes they are in ours. They are people (Edmund and Lucy had once met one). So now they found showers of glittering people, all with long hair like burning silver and spears like white-hot metal, rushing down to them out of the black air, swifter than falling stones. They made a hissing noise as they landed and burnt the grass. And all these stars glided past them and stood somewhere behind, a little to the right.'

So not quite divinities, perhaps, but not far off. A bit like the Sidhe of Gaelic mythology, maybe.

And in Lewis's Space Trilogy the planets are most certainly alive, especially in the last volume, That Hideous Strength, where they each play a decisive role in the affairs of the Earth.


\hfill

\end{sffamily}
\end{footnotesize}

